Nous avons développé des preuves de concept de visualisations pour des images à caractère scientifique, telles que des diagrammes astronomiques ou des figures géométriques. Toutefois, ces approches peuvent être appliquées à tout corpus cherchant à étudier des motifs similaires entre différentes images, afin d'en retracer la diffusion dans le temps et l'espace.


\subsection{Les photographies de presse}

Le projet \textit{High Vision}\footcite{HIGHVISIONProjet2025} étudie la diffusion des photographies de presse de la fin du \XIX siècle aux années 1940. Ce corpus est particulièrement pertinent car différents journaux réutilisaient les mêmes photographies d'agence de presse. Une des pratiques courantes était le photo-montage. Les journaux réutilisaient les mêmes photographies mais en les modifiant de manière différente. Le but du projet est donc de retracer la diffusion d'une même photographie dans différents journaux. Le projet \textit{High Vision} a exprimé sa volonté d'utiliser l'application AIKON. Il pourra donc bénéficier des visualisations quand ces dernières seront intégrées à l'application.

\subsection{Les peintures}

Les peintures constituent également un type de corpus adapté aux visualisations que nous avons réalisé. Nous pouvons prendre comme exemple le \textit{Brueghel dataset} utilisé dans le cadre du projet \gls{enherit}\footcite{EnHeritEnhancingHeritage}. Dans le cadre des ateliers de peintres comme celui de Brueghel, il n'était pas rare de réutiliser plusieurs fois des motifs créés dans des croquis préliminaires ou des études en les modifiant ou non\footcite{shenDiscoveringVisualPatterns2019}. Les visualisations que nous avons produites peuvent donc servir à établir les liens entre ces différentes œuvres pour retracer la diffusion d'un motif.

\subsection{Les objets archéologiques}

Les visualisations que nous avons produites peuvent s'appliquer sur des corpus différents de ceux en deux dimensions, notamment dans le domaine de l'archéologie. En effet, il est courant que des objets sculptés réutilisent les mêmes motifs. 

Récemment, des chercheurs ont publié un article présentant leurs travaux sur la diffusion des motifs sur des tuiles chinoises sur une période de 2000 ans. Ces derniers utilisent l'intelligence artificielle à la fois pour restaurer les tuiles dégradées par le temps mais aussi pour analyser leurs similarités. Ils ont d'ailleurs développé un prototype proposant une visualisation similaire aux nôtres\footcite{yuEmbodiedHeritageIntegrating2025a}. 